\begin{abstract}
Multi-teacher on-policy distillation transfers domain expertise into a shared student. We study how loss averaging, optimizer history, and vocabulary supervision shape this transfer, using Qwen3-1.7B and four domain specialists, with additional SmolLM3-3B measurements. Token and response averaging produce different gradients whose update directions become more similar after applying saved Adam moments; final task scores are close under top-64 supervision. Training changes about 97\% of FP32 parameters, with only 7--9\% remaining visibly changed after BF16 rounding. Adam's accumulated moments also account for much of the alignment between teacher-specific updates: subtracting the zero-gradient step reduces their mean cosine from above 0.83 to below 0.01 on domain-specific inputs. Expanding supervision from a sampled token to the top-64 intersection improves final mathematics accuracy and lowers science accuracy, with relative gains changing during training. Top-64 closely approximates the full-vocabulary gradient direction in the high-coverage Qwen measurements; SmolLM measurements with lower coverage retain appreciable gradient error. Together, these findings connect local update measurements to task-dependent learning trajectories and identify optimizer history and numerical precision as central to interpreting distillation updates.
\end{abstract}
